\section{Extensions to imperfect products and natural probes}
\label{app:enhancement-theory}

\subsection{Joint product and response perturbations}

The same-product restriction in Theorem~\ref{thm:response_stability} can be
relaxed for small product perturbations by constructing an intermediate factorization with exactly
the reference product. This is an inverse-property result for valid factors;
it does not certify the unconstrained numerical recovery algorithm.

\begin{corollary}[Joint product--response stability]
\label{cor:joint-stability}
Let two factor pairs satisfy all hypotheses of
Theorem~\ref{thm:response_stability} except equal products, with the same
constants $\alpha,s_A,s_B,L_A,L_B$ and known positive rates and temperature.
Write
\[
 \epsilon_\Theta=\|AB-A'B'\|_F,\qquad
 \epsilon_R=\|\mathcal R_{A,B}-\mathcal R_{A',B'}\|_{F\to F}.
\]
Assume $\epsilon_\Theta\le\min\{\alpha s_B/2,s_A s_B/4\}$.
Let $C_*$ be the explicit constant of that theorem evaluated at
$(\alpha/2,s_A/2,s_B/2,2L_A,2L_B)$, and put
\[
 \kappa=(s_B^{-2}+4s_A^{-2})^{1/2},\quad M=2L_B,\quad
 L_R=\bigl[(2\eta_B\sqrt L+6\lambda_Z M^2)^2+(2\lambda_Z M)^2\bigr]^{1/2}.
\]
Then the joint Frobenius distance modulo a common permutation satisfies
\[
 d_{\rm perm}((A,B),(A',B'))
 \le C_*\epsilon_R+\kappa(1+C_*L_R)\epsilon_\Theta.
\]
\end{corollary}

\begin{proof}
Let $P_A=AA^\dagger$, $\widetilde A=P_AA'$, and
$\widetilde B=\widetilde A^\dagger AB$. Since $A\ones=\ones$,
$P_A\ones=\ones$ and $\widetilde A\ones=\ones$. Moreover,
\[
 (I-P_A)A'=(I-P_A)(A'B'-AB)B'^\dagger,
 \qquad \|\widetilde A-A'\|_F\le\epsilon_\Theta/s_B.
\]
Thus $\min\widetilde A\ge\alpha/2$ and
$\sigma_{\min}(\widetilde A)\ge s_A/2$. Its column space is that of $A$, so
$\widetilde A\widetilde B=AB$. Also
\[
 \widetilde B-B'=\widetilde A^\dagger P_A(AB-A'B'),\qquad
 \|\widetilde B-B'\|_F\le2\epsilon_\Theta/s_A.
\]
Consequently $\sigma_{\min}(\widetilde B)\ge s_B/2$, and both pairs
$(A,B),(\widetilde A,\widetilde B)$ obey the enlarged domain bounds defining
$C_*$. The joint displacement from $(A',B')$ to the intermediate pair is at
most $\kappa\epsilon_\Theta$.

For simplex rows $a,a'$, let $C(a)=\operatorname{diag}(a)-aa^\top$.
The elementary bounds $\|a\|_2\le1$, $\|C(a)\|_2\le1$, and
$\|C(a)-C(a')\|_2\le3\|a-a'\|_2$ imply
$\|C(a)^2-C(a')^2\|_2\le6\|a-a'\|_2$.
In the response formula, the Gram contribution changes by at most
$2\eta_B\sqrt L\|\Delta A\|_F$. Each local diagonal block changes by at most
$\lambda_Z(6M^2\|\Delta a_i\|_2+2M\|\Delta B\|_F)$.
The norm of the direct sum is the maximum block norm; therefore the full
response is $L_R$-Lipschitz in the joint factor norm on this domain.
The response discrepancy between $(A,B)$ and the intermediate pair is bounded
by $\epsilon_R+L_R\kappa\epsilon_\Theta$. Apply the same-product theorem and
then the triangle inequality; common permutations preserve the joint norm.
\end{proof}

If two valid feasible pairs each match the same noisy observations within
$e_\Theta,e_R$, the corollary applies with $2e_\Theta,2e_R$, provided the
small-product-error condition holds. It does not establish existence of a
feasible pair, an efficient constrained solver, or stability near rank loss.
The constants are conservative and not numerical accuracy predictions.

\subsection{A sufficient condition for tomography with natural gradients}
\label{app:natural-tomography}

Let $U\in\RR^{K\times D}$ have orthonormal rows spanning the exact product's
row space. Write the structured response as
$R^{(p)}=QG^{(p)}+\mathcal T(G^{(p)}U^\top)U$, where
$Q=\eta_B AA^\top$ and $\mathcal T$ applies a $K\times K$ matrix $T_i$ to each
row coordinate, using column-vector convention. Define
$G_\perp^{(p)}=G^{(p)}-(G^{(p)}U^\top)U$ and similarly $R_\perp^{(p)}$.

\begin{proposition}[Rank-checked natural-probe reconstruction]
\label{prop:natural-probes}
If $H=\sum_p G_\perp^{(p)}G_\perp^{(p)\top}$ is invertible and, for every
layer $i$, the vectors $(G^{(p)}U^\top)_{i,:}$ span $\RR^K$, the observations
uniquely determine $Q$ and all $T_i$. They therefore determine the response
on all gradients within the assumed response family.
\end{proposition}
\begin{proof}
Projection eliminates the local terms, giving
$R_\perp^{(p)}=QG_\perp^{(p)}$. Hence
\[
 Q=\left(\sum_pR_\perp^{(p)}G_\perp^{(p)\top}\right)H^{-1}.
\]
Subtract $QG^{(p)}$ and project onto $U$; each layer's remaining linear
system has a unique $T_i$ by its spanning assumption.
\end{proof}
This is a sufficient algebraic condition, not a guarantee about task gradients
or an optimal-query result. The implementation checks numerical ranks and
conditioning, and uses fixed pseudoinverse tolerances. Small singular values
can amplify measurement noise. The designed probes avoid relying on empirical
task-gradient richness; natural probes trade that design for a condition that
must be checked at the checkpoint. Measurements here still use controlled
Euclidean responses at a fixed checkpoint, not an observed AdamW trajectory.
